% ============================================================
\chapter{Fonctions Convexes}
\label{ch:fonctions}
% ============================================================

Si les ensembles convexes sont le d\'ecor, les fonctions convexes en sont
les protagonistes. Une fonction est convexe si elle \og courbe vers le
haut\fg~: la s\'ecante entre deux points de son graphe est toujours
au-dessus du graphe. Cette propri\'et\'e, d'apparence modeste, a des
cons\'equences spectaculaires~: tout minimum local est global, les
conditions d'optimalit\'e du premier ordre sont suffisantes, et les
algorithmes d'optimisation convergent avec des garanties th\'eoriques
solides. Rockafellar a dit un jour que les fonctions convexes sont aux
in\'egalit\'es ce que les fonctions affines sont aux \'egalit\'es --- une
analogie profonde que ce chapitre va explorer.

% ------------------------------------------------------------
\section{Introduction}
% ------------------------------------------------------------

Les fonctions convexes sont au c\oe ur de l'optimisation convexe. Leur propri\'et\'e
fondamentale --- tout minimum local est global --- rend les probl\`emes d'optimisation
convexe \og{}bien pos\'es\fg et th\'eoriquement tractables. Ce chapitre pr\'esente
les d\'efinitions, caract\'erisations et propri\'et\'es essentielles.

\section{D\'efinitions et premi\`eres propri\'et\'es}

\begin{definition}[Fonction convexe]
Une fonction $f : \R^n \to \R \cup \{+\infty\}$ est \emph{convexe} si son domaine
effectif $\mathrm{dom}(f) = \{x \mid f(x) < +\infty\}$ est convexe et si pour tout
$x, y \in \mathrm{dom}(f)$ et tout $\theta \in [0,1]$~:
\[
    f(\theta x + (1-\theta)y) \leq \theta f(x) + (1-\theta)f(y).
\]
\end{definition}

\begin{definition}[Convexit\'e stricte]
$f$ est \emph{strictement convexe} si l'in\'egalit\'e ci-dessus est stricte pour
$x \neq y$ et $\theta \in (0,1)$.
\end{definition}

\begin{definition}[Forte convexit\'e]
$f$ est \emph{fortement convexe} de param\`etre $m > 0$ (ou $m$-fortement convexe) si
$f(x) - \frac{m}{2}\norm{x}^2$ est convexe, c'est-\`a-dire pour tout $x, y$ et
$\theta \in [0,1]$~:
\[
    f(\theta x + (1-\theta)y) \leq \theta f(x) + (1-\theta)f(y)
    - \frac{m}{2}\theta(1-\theta)\norm{x-y}^2.
\]
\end{definition}

\begin{formules}[title=\textbf{Hi\'erarchie de convexit\'e}]
\[
    \text{fortement convexe} \;\Rightarrow\; \text{strictement convexe}
    \;\Rightarrow\; \text{convexe}.
\]
Les implications inverses sont fausses en g\'en\'eral.
\end{formules}

\section{\'Epigraphe et caract\'erisation}

\begin{definition}[\'Epigraphe]
L'\emph{\'epigraphe} d'une fonction $f : \R^n \to \R \cup \{+\infty\}$ est~:
\[
    \mathrm{epi}(f) = \{(x, t) \in \R^n \times \R \mid f(x) \leq t\}.
\]
\end{definition}

\begin{theoreme}[Caract\'erisation \'epigraphique]
\label{thm:epigraphe}
$f$ est convexe si et seulement si $\mathrm{epi}(f)$ est un ensemble convexe de
$\R^{n+1}$.
\end{theoreme}

\begin{proof}
$(\Rightarrow)$ Soient $(x_1, t_1), (x_2, t_2) \in \mathrm{epi}(f)$ et
$\theta \in [0,1]$. Alors~:
\[
    f(\theta x_1 + (1-\theta)x_2) \leq \theta f(x_1) + (1-\theta)f(x_2)
    \leq \theta t_1 + (1-\theta)t_2,
\]
donc $(\theta x_1 + (1-\theta)x_2, \theta t_1 + (1-\theta)t_2) \in \mathrm{epi}(f)$.

$(\Leftarrow)$ Soient $x_1, x_2 \in \mathrm{dom}(f)$. Alors
$(x_1, f(x_1)), (x_2, f(x_2)) \in \mathrm{epi}(f)$. Par convexit\'e de l'\'epigraphe~:
\[
    (\theta x_1 + (1-\theta)x_2, \theta f(x_1) + (1-\theta)f(x_2)) \in \mathrm{epi}(f),
\]
ce qui signifie $f(\theta x_1 + (1-\theta)x_2) \leq \theta f(x_1) + (1-\theta)f(x_2)$.
\end{proof}

\section{Ensembles de sous-niveau}

\begin{definition}[Ensemble de sous-niveau]
L'\emph{ensemble de sous-niveau $\alpha$} de $f$ est~:
\[
    S_\alpha = \{x \in \mathrm{dom}(f) \mid f(x) \leq \alpha\}.
\]
\end{definition}

\begin{proposition}
Si $f$ est convexe, alors tous ses ensembles de sous-niveau sont convexes.
\end{proposition}

\begin{attention}[title=\textbf{R\'eciproque fausse}]
La r\'eciproque est fausse~: $f(x) = e^x$ a des sous-niveaux convexes (intervalles)
mais $g(x) = \sqrt{\abs{x}}$ a aussi des sous-niveaux convexes sans \^etre convexe.
Une fonction dont les sous-niveaux sont convexes est dite \emph{quasi-convexe}.
\end{attention}

\section{Caract\'erisations diff\'erentielles}

\subsection{Condition du premier ordre}

\begin{theoreme}[Condition du premier ordre]
\label{thm:premier_ordre}
Soit $f : \R^n \to \R$ diff\'erentiable sur un ouvert convexe $\Omega$. Alors
$f$ est convexe sur $\Omega$ si et seulement si~:
\[
    f(y) \geq f(x) + \inner{\nabla f(x)}{y - x}, \quad \forall x, y \in \Omega.
\]
\end{theoreme}

\begin{proof}
$(\Rightarrow)$ Par convexit\'e, pour $t \in (0,1]$~:
\[
    f(x + t(y-x)) \leq f(x) + t(f(y) - f(x)).
\]
Donc $\frac{f(x+t(y-x)) - f(x)}{t} \leq f(y) - f(x)$. En passant \`a la limite
$t \to 0^+$~: $\inner{\nabla f(x)}{y-x} \leq f(y) - f(x)$.

$(\Leftarrow)$ Soient $x, y \in \Omega$ et $z = \theta x + (1-\theta)y$. Alors~:
\begin{align*}
    f(x) &\geq f(z) + \inner{\nabla f(z)}{x - z}, \\
    f(y) &\geq f(z) + \inner{\nabla f(z)}{y - z}.
\end{align*}
En multipliant par $\theta$ et $1-\theta$ respectivement et en sommant~:
\[
    \theta f(x) + (1-\theta)f(y) \geq f(z) + \inner{\nabla f(z)}{\theta x + (1-\theta)y - z} = f(z).
\]
\end{proof}

\begin{intuition}
La condition du premier ordre signifie que le graphe de $f$ est toujours au-dessus
de ses tangentes. G\'eom\'etriquement, la surface $z = f(x)$ est \og{}bomb\'ee vers
le haut\fg.
\end{intuition}

\subsection{Condition du second ordre}

\begin{theoreme}[Condition du second ordre]
\label{thm:second_ordre}
Soit $f : \R^n \to \R$ deux fois diff\'erentiable sur un ouvert convexe $\Omega$. Alors~:
\begin{enumerate}
    \item $f$ est convexe sur $\Omega$ $\Leftrightarrow$ $\nabla^2 f(x) \succeq 0$
          pour tout $x \in \Omega$.
    \item $f$ est fortement convexe de param\`etre $m$ $\Leftrightarrow$
          $\nabla^2 f(x) \succeq mI$ pour tout $x \in \Omega$.
    \item Si $\nabla^2 f(x) \succ 0$ pour tout $x$, alors $f$ est strictement convexe
          (la r\'eciproque est fausse).
\end{enumerate}
\end{theoreme}

\begin{proof}[Preuve de (1), $\Rightarrow$]
Par la condition du premier ordre, pour tout $h \in \R^n$ et $t > 0$ assez petit~:
\[
    f(x + th) \geq f(x) + t\inner{\nabla f(x)}{h}.
\]
Par d\'eveloppement de Taylor~:
\[
    f(x+th) = f(x) + t\inner{\nabla f(x)}{h} + \frac{t^2}{2}h^T\nabla^2 f(x)h + o(t^2).
\]
Donc $h^T\nabla^2 f(x)h + o(t^2)/t^2 \geq 0$. En faisant $t \to 0$~:
$h^T\nabla^2 f(x)h \geq 0$.
\end{proof}

\section{In\'egalit\'e de Jensen}

\begin{theoreme}[In\'egalit\'e de Jensen]
Soit $f$ convexe et $x_1, \ldots, x_k \in \mathrm{dom}(f)$ avec
$\theta_1, \ldots, \theta_k \geq 0$, $\sum \theta_i = 1$. Alors~:
\[
    f\left(\sum_{i=1}^k \theta_i x_i\right) \leq \sum_{i=1}^k \theta_i f(x_i).
\]
\end{theoreme}

\begin{corollaire}[Jensen probabiliste]
Soit $X$ une variable al\'eatoire \`a valeurs dans $\mathrm{dom}(f)$ avec
$\mathbb{E}[\abs{X}] < \infty$. Si $f$ est convexe~:
\[
    f(\mathbb{E}[X]) \leq \mathbb{E}[f(X)].
\]
\end{corollaire}

\section{Op\'erations pr\'eservant la convexit\'e des fonctions}

\begin{theoreme}[Op\'erations sur les fonctions convexes]
\label{thm:operations_fonctions}
\begin{enumerate}
    \item \textbf{Combinaison positive}~: si $f_1, \ldots, f_k$ sont convexes et
          $\alpha_1, \ldots, \alpha_k \geq 0$, alors $\sum \alpha_i f_i$ est convexe.
    \item \textbf{Maximum ponctuel}~: si $f_1, \ldots, f_k$ sont convexes, alors
          $g(x) = \max_i f_i(x)$ est convexe.
    \item \textbf{Supremum}~: si $(f_\alpha)_{\alpha \in I}$ est une famille de
          fonctions convexes, alors $g(x) = \sup_\alpha f_\alpha(x)$ est convexe.
    \item \textbf{Composition affine}~: si $f$ est convexe, alors $g(x) = f(Ax + b)$
          est convexe.
    \item \textbf{Infimale convolution}~: si $f, g$ sont convexes, alors
          $(f \mathop{\square} g)(x) = \inf_y \{f(y) + g(x-y)\}$ est convexe.
\end{enumerate}
\end{theoreme}

\section{Convexit\'e et lissit\'e}

\begin{definition}[$L$-lissit\'e]
Une fonction convexe diff\'erentiable $f$ est \emph{$L$-lisse} si son gradient est
$L$-lipschitzien~:
\[
    \norm{\nabla f(x) - \nabla f(y)} \leq L\norm{x - y}, \quad \forall x, y.
\]
\end{definition}

\begin{theoreme}[Caract\'erisations de la $L$-lissit\'e]
\label{thm:lissage}
Pour $f$ convexe et diff\'erentiable, les conditions suivantes sont \'equivalentes~:
\begin{enumerate}
    \item $\nabla f$ est $L$-lipschitzien.
    \item $f(y) \leq f(x) + \inner{\nabla f(x)}{y-x} + \frac{L}{2}\norm{y-x}^2$ pour tout $x, y$.
    \item $f(y) \geq f(x) + \inner{\nabla f(x)}{y-x} + \frac{1}{2L}\norm{\nabla f(y) - \nabla f(x)}^2$ pour tout $x, y$.
    \item $\inner{\nabla f(x) - \nabla f(y)}{x - y} \geq \frac{1}{L}\norm{\nabla f(x) - \nabla f(y)}^2$ pour tout $x, y$.
\end{enumerate}
Si $f$ est deux fois diff\'erentiable, ces conditions \'equivalent \`a $\nabla^2 f(x) \preceq LI$.
\end{theoreme}

\begin{formules}[title=\textbf{Encadrement fondamental}]
Pour $f$ convexe, $m$-fortement convexe et $L$-lisse~:
\[
    f(x) + \inner{\nabla f(x)}{y-x} + \frac{m}{2}\norm{y-x}^2
    \leq f(y) \leq
    f(x) + \inner{\nabla f(x)}{y-x} + \frac{L}{2}\norm{y-x}^2.
\]
Le nombre de conditionnement $\kappa = L/m$ mesure la difficult\'e du probl\`eme.
\end{formules}

\section{Exemples importants}

\begin{exemple}[Fonctions convexes classiques]
\begin{enumerate}
    \item \textbf{Fonctions affines}~: $f(x) = \inner{a}{x} + b$ (convexe et concave).
    \item \textbf{Normes}~: $\norm{\cdot}_p$ pour $p \geq 1$ (convexe).
    \item \textbf{Quadratique}~: $f(x) = \frac{1}{2}x^TQx + q^Tx + c$ avec $Q \succeq 0$.
          Fortement convexe ssi $Q \succ 0$, avec $m = \lambda_{\min}(Q)$, $L = \lambda_{\max}(Q)$.
    \item \textbf{Exponentielle}~: $f(x) = e^{ax}$ (convexe sur $\R$).
    \item \textbf{Log-sum-exp}~: $f(x) = \log(\sum_i e^{x_i})$ (convexe sur $\R^n$).
    \item \textbf{Entropie n\'egative}~: $f(x) = \sum_i x_i \log x_i$ sur $\R^n_{++}$.
    \item \textbf{Perte logistique}~: $f(x) = \log(1 + e^{-x})$ (convexe, lisse).
\end{enumerate}
\end{exemple}

\section{Continuit\'e des fonctions convexes}

\begin{theoreme}[Continuit\'e]
\label{thm:continuite}
Toute fonction convexe $f : \R^n \to \R$ (finie partout) est continue. Plus
pr\'ecis\'ement, toute fonction convexe propre est continue sur l'int\'erieur
de son domaine effectif.
\end{theoreme}

\begin{theoreme}[Diff\'erentiabilit\'e presque partout]
Toute fonction convexe $f : \R^n \to \R$ est diff\'erentiable presque partout
(au sens de Lebesgue). En dimension $1$, $f$ admet des d\'eriv\'ees \`a gauche
et \`a droite en tout point.
\end{theoreme}

\section{Fonctions convexes conjugu\'ees --- aper\c{c}u}

\begin{definition}[Conjugu\'ee de Fenchel]
La \emph{conjugu\'ee de Fenchel} (ou transform\'ee de Legendre--Fenchel) de
$f : \R^n \to \R \cup \{+\infty\}$ est~:
\[
    f^*(y) = \sup_{x \in \R^n} \left\{\inner{y}{x} - f(x)\right\}.
\]
\end{definition}

\begin{remarque}
$f^*$ est toujours convexe et semi-continue inf\'erieurement (comme supremum de
fonctions affines en $y$), m\^eme si $f$ ne l'est pas. La conjugaison convexe sera
\'etudi\'ee en d\'etail au chapitre~\ref{ch:fenchel}.
\end{remarque}

\section{Impl\'ementation Python}

\begin{codeblock}[title=\textbf{V\'erification de convexit\'e et visualisation}]
\begin{minted}{python}
import numpy as np
import matplotlib.pyplot as plt

def check_convexity_numerical(f, x_range, n_tests=1000):
    """Verification numerique de convexite en 1D."""
    for _ in range(n_tests):
        x = np.random.uniform(*x_range)
        y = np.random.uniform(*x_range)
        t = np.random.uniform(0, 1)
        lhs = f(t * x + (1 - t) * y)
        rhs = t * f(x) + (1 - t) * f(y)
        if lhs > rhs + 1e-10:
            return False
    return True

# Exemples
functions = {
    'x^2': lambda x: x**2,
    'exp(x)': lambda x: np.exp(x),
    '|x|': lambda x: np.abs(x),
    'log(1+exp(x))': lambda x: np.log(1 + np.exp(x)),
}

fig, axes = plt.subplots(2, 2, figsize=(10, 8))
for ax, (name, f) in zip(axes.flat, functions.items()):
    x = np.linspace(-3, 3, 200)
    ax.plot(x, f(x), 'b-', linewidth=2)
    # Tangente en x=1
    x0 = 1.0
    h = 1e-5
    grad = (f(x0 + h) - f(x0 - h)) / (2 * h)
    tangent = f(x0) + grad * (x - x0)
    ax.plot(x, tangent, 'r--', alpha=0.7)
    is_cvx = check_convexity_numerical(f, (-3, 3))
    ax.set_title(f'{name} (convexe: {is_cvx})')
    ax.grid(True, alpha=0.3)
plt.tight_layout()
plt.savefig('convex_functions.pdf')
\end{minted}
\end{codeblock}

\section{Exercices}

\begin{exercice}[$\star$]
Montrer que $f(x) = \max(x_1, \ldots, x_n)$ est convexe sur $\R^n$.
\end{exercice}

\begin{exercice}[$\star$]
Montrer que si $f$ est convexe et $g$ est affine, alors $f \circ g$ est convexe.
\end{exercice}

\begin{exercice}[$\star\star$]
Montrer qu'une fonction $f : \R^n \to \R$ diff\'erentiable est $m$-fortement convexe
si et seulement si $\inner{\nabla f(x) - \nabla f(y)}{x - y} \geq m\norm{x - y}^2$
pour tout $x, y$.
\end{exercice}

\begin{exercice}[$\star\star$]
Montrer que $f(x) = \log\left(\sum_{i=1}^n e^{x_i}\right)$ est convexe et calculer
sa constante de lissit\'e $L$ par rapport \`a la norme $\ell_2$.
\end{exercice}

\begin{exercice}[$\star\star$]
Soit $f$ convexe, propre et semi-continue inf\'erieurement. Montrer que $f$ est
minor\'ee par une fonction affine.
\end{exercice}

\begin{exercice}[$\star\star\star$]
Montrer l'\'equivalence des quatre caract\'erisations de la $L$-lissit\'e du
th\'eor\`eme~\ref{thm:lissage}.
\end{exercice}

\begin{exercice}[$\star\star\star$]
Soit $f$ convexe sur $\R^n$ et $x^*$ un minimiseur. Montrer que si $f$ est
$m$-fortement convexe~:
\[
    f(x) - f(x^*) \geq \frac{m}{2}\norm{x - x^*}^2, \quad \forall x \in \R^n.
\]
\end{exercice}
